\IfFileExists{stacks-project.cls}{%
\documentclass{stacks-project}
}{%
\documentclass{amsart}
}

% For dealing with references we use the comment environment
\usepackage{verbatim}
\newenvironment{reference}{\comment}{\endcomment}
%\newenvironment{reference}{}{}
\newenvironment{slogan}{\comment}{\endcomment}
\newenvironment{history}{\comment}{\endcomment}

% For commutative diagrams we use Xy-pic
\usepackage[all]{xy}

% We use 2cell for 2-commutative diagrams.
\xyoption{2cell}
\UseAllTwocells

% We use multicol for the list of chapters between chapters
\usepackage{multicol}

% This is generally recommended for better output
\usepackage{lmodern}
\usepackage[T1]{fontenc}

% For cross-file-references
\usepackage{xr-hyper}

% Package for hypertext links:
\usepackage{hyperref}

% For any local file, say "hello.tex" you want to link to please
% use \externaldocument[hello-]{hello}
\externaldocument[introduction-]{introduction}
\externaldocument[conventions-]{conventions}
\externaldocument[sets-]{sets}
\externaldocument[categories-]{categories}
\externaldocument[topology-]{topology}
\externaldocument[sheaves-]{sheaves}
\externaldocument[sites-]{sites}
\externaldocument[stacks-]{stacks}
\externaldocument[fields-]{fields}
\externaldocument[algebra-]{algebra}
\externaldocument[brauer-]{brauer}
\externaldocument[homology-]{homology}
\externaldocument[derived-]{derived}
\externaldocument[simplicial-]{simplicial}
\externaldocument[more-algebra-]{more-algebra}
\externaldocument[smoothing-]{smoothing}
\externaldocument[modules-]{modules}
\externaldocument[sites-modules-]{sites-modules}
\externaldocument[injectives-]{injectives}
\externaldocument[cohomology-]{cohomology}
\externaldocument[sites-cohomology-]{sites-cohomology}
\externaldocument[dga-]{dga}
\externaldocument[dpa-]{dpa}
\externaldocument[sdga-]{sdga}
\externaldocument[hypercovering-]{hypercovering}
\externaldocument[schemes-]{schemes}
\externaldocument[constructions-]{constructions}
\externaldocument[properties-]{properties}
\externaldocument[morphisms-]{morphisms}
\externaldocument[coherent-]{coherent}
\externaldocument[divisors-]{divisors}
\externaldocument[limits-]{limits}
\externaldocument[varieties-]{varieties}
\externaldocument[topologies-]{topologies}
\externaldocument[descent-]{descent}
\externaldocument[perfect-]{perfect}
\externaldocument[more-morphisms-]{more-morphisms}
\externaldocument[flat-]{flat}
\externaldocument[groupoids-]{groupoids}
\externaldocument[more-groupoids-]{more-groupoids}
\externaldocument[etale-]{etale}
\externaldocument[chow-]{chow}
\externaldocument[intersection-]{intersection}
\externaldocument[pic-]{pic}
\externaldocument[weil-]{weil}
\externaldocument[adequate-]{adequate}
\externaldocument[dualizing-]{dualizing}
\externaldocument[duality-]{duality}
\externaldocument[discriminant-]{discriminant}
\externaldocument[derham-]{derham}
\externaldocument[local-cohomology-]{local-cohomology}
\externaldocument[algebraization-]{algebraization}
\externaldocument[curves-]{curves}
\externaldocument[resolve-]{resolve}
\externaldocument[models-]{models}
\externaldocument[functors-]{functors}
\externaldocument[equiv-]{equiv}
\externaldocument[pione-]{pione}
\externaldocument[etale-cohomology-]{etale-cohomology}
\externaldocument[proetale-]{proetale}
\externaldocument[relative-cycles-]{relative-cycles}
\externaldocument[more-etale-]{more-etale}
\externaldocument[trace-]{trace}
\externaldocument[crystalline-]{crystalline}
\externaldocument[spaces-]{spaces}
\externaldocument[spaces-properties-]{spaces-properties}
\externaldocument[spaces-morphisms-]{spaces-morphisms}
\externaldocument[decent-spaces-]{decent-spaces}
\externaldocument[spaces-cohomology-]{spaces-cohomology}
\externaldocument[spaces-limits-]{spaces-limits}
\externaldocument[spaces-divisors-]{spaces-divisors}
\externaldocument[spaces-over-fields-]{spaces-over-fields}
\externaldocument[spaces-topologies-]{spaces-topologies}
\externaldocument[spaces-descent-]{spaces-descent}
\externaldocument[spaces-perfect-]{spaces-perfect}
\externaldocument[spaces-more-morphisms-]{spaces-more-morphisms}
\externaldocument[spaces-flat-]{spaces-flat}
\externaldocument[spaces-groupoids-]{spaces-groupoids}
\externaldocument[spaces-more-groupoids-]{spaces-more-groupoids}
\externaldocument[bootstrap-]{bootstrap}
\externaldocument[spaces-pushouts-]{spaces-pushouts}
\externaldocument[spaces-chow-]{spaces-chow}
\externaldocument[groupoids-quotients-]{groupoids-quotients}
\externaldocument[spaces-more-cohomology-]{spaces-more-cohomology}
\externaldocument[spaces-simplicial-]{spaces-simplicial}
\externaldocument[spaces-duality-]{spaces-duality}
\externaldocument[formal-spaces-]{formal-spaces}
\externaldocument[restricted-]{restricted}
\externaldocument[spaces-resolve-]{spaces-resolve}
\externaldocument[formal-defos-]{formal-defos}
\externaldocument[defos-]{defos}
\externaldocument[cotangent-]{cotangent}
\externaldocument[examples-defos-]{examples-defos}
\externaldocument[algebraic-]{algebraic}
\externaldocument[examples-stacks-]{examples-stacks}
\externaldocument[stacks-sheaves-]{stacks-sheaves}
\externaldocument[criteria-]{criteria}
\externaldocument[artin-]{artin}
\externaldocument[quot-]{quot}
\externaldocument[stacks-properties-]{stacks-properties}
\externaldocument[stacks-morphisms-]{stacks-morphisms}
\externaldocument[stacks-limits-]{stacks-limits}
\externaldocument[stacks-cohomology-]{stacks-cohomology}
\externaldocument[stacks-perfect-]{stacks-perfect}
\externaldocument[stacks-introduction-]{stacks-introduction}
\externaldocument[stacks-more-morphisms-]{stacks-more-morphisms}
\externaldocument[stacks-geometry-]{stacks-geometry}
\externaldocument[moduli-]{moduli}
\externaldocument[moduli-curves-]{moduli-curves}
\externaldocument[examples-]{examples}
\externaldocument[exercises-]{exercises}
\externaldocument[guide-]{guide}
\externaldocument[desirables-]{desirables}
\externaldocument[coding-]{coding}
\externaldocument[obsolete-]{obsolete}
\externaldocument[fdl-]{fdl}
\externaldocument[index-]{index}

% Theorem environments.
%
\theoremstyle{plain}
\newtheorem{theorem}[subsection]{Theorem}
\newtheorem{proposition}[subsection]{Proposition}
\newtheorem{lemma}[subsection]{Lemma}

\theoremstyle{definition}
\newtheorem{definition}[subsection]{Definition}
\newtheorem{example}[subsection]{Example}
\newtheorem{exercise}[subsection]{Exercise}
\newtheorem{situation}[subsection]{Situation}

\theoremstyle{remark}
\newtheorem{remark}[subsection]{Remark}
\newtheorem{remarks}[subsection]{Remarks}

\numberwithin{equation}{subsection}

% Macros
%
\def\lim{\mathop{\mathrm{lim}}\nolimits}
\def\colim{\mathop{\mathrm{colim}}\nolimits}
\def\Spec{\mathop{\mathrm{Spec}}}
\def\Hom{\mathop{\mathrm{Hom}}\nolimits}
\def\Ext{\mathop{\mathrm{Ext}}\nolimits}
\def\SheafHom{\mathop{\mathcal{H}\!\mathit{om}}\nolimits}
\def\SheafExt{\mathop{\mathcal{E}\!\mathit{xt}}\nolimits}
\def\Sch{\mathit{Sch}}
\def\Mor{\mathop{\mathrm{Mor}}\nolimits}
\def\Ob{\mathop{\mathrm{Ob}}\nolimits}
\def\Sh{\mathop{\mathit{Sh}}\nolimits}
\def\NL{\mathop{N\!L}\nolimits}
\def\CH{\mathop{\mathrm{CH}}\nolimits}
\def\proetale{{pro\text{-}\acute{e}tale}}
\def\etale{{\acute{e}tale}}
\def\QCoh{\mathit{QCoh}}
\def\Ker{\mathop{\mathrm{Ker}}}
\def\Im{\mathop{\mathrm{Im}}}
\def\Coker{\mathop{\mathrm{Coker}}}
\def\Coim{\mathop{\mathrm{Coim}}}

% Boxtimes
%
\DeclareMathSymbol{\boxtimes}{\mathbin}{AMSa}{"02}

%
% Macros for moduli stacks/spaces
%
\def\QCohstack{\mathcal{QC}\!\mathit{oh}}
\def\Cohstack{\mathcal{C}\!\mathit{oh}}
\def\Spacesstack{\mathcal{S}\!\mathit{paces}}
\def\Quotfunctor{\mathrm{Quot}}
\def\Hilbfunctor{\mathrm{Hilb}}
\def\Curvesstack{\mathcal{C}\!\mathit{urves}}
\def\Polarizedstack{\mathcal{P}\!\mathit{olarized}}
\def\Complexesstack{\mathcal{C}\!\mathit{omplexes}}
% \Pic is the operator that assigns to X its picard group, usage \Pic(X)
% \Picardstack_{X/B} denotes the Picard stack of X over B
% \Picardfunctor_{X/B} denotes the Picard functor of X over B
\def\Pic{\mathop{\mathrm{Pic}}\nolimits}
\def\Picardstack{\mathcal{P}\!\mathit{ic}}
\def\Picardfunctor{\mathrm{Pic}}
\def\Deformationcategory{\mathcal{D}\!\mathit{ef}}
\usepackage{cleveref}
\crefname{originallemma}{original lemma}{original lemmas}
\Crefname{originallemma}{Original lemma}{Original lemmas}
\makeatletter
\newlabel{lemma-stalk-R-f-shriek}{{0GK0}{source}{}{sourcepointer-1}{}}
\newlabel{lemma-stalk-R-f-shriek@cref}{{[originallemma][0][]0GK0}{[1][0][]source}{}{}{}}
\makeatother
\crefname{originallemma}{original lemma}{original lemmas}
\Crefname{originallemma}{Original lemma}{Original lemmas}
\makeatletter
\newlabel{lemma-finiteness-curves}{{0GKR}{source}{}{sourcepointer-2}{}}
\newlabel{lemma-finiteness-curves@cref}{{[originallemma][0][]0GKR}{[1][0][]source}{}{}{}}
\makeatother
\crefname{originallemma}{original lemma}{original lemmas}
\Crefname{originallemma}{Original lemma}{Original lemmas}
\makeatletter
\newlabel{lemma-base-change-shriek}{{0F7L}{source}{}{sourcepointer-3}{}}
\newlabel{lemma-base-change-shriek@cref}{{[originallemma][0][]0F7L}{[1][0][]source}{}{}{}}
\makeatother
\crefname{originallemma}{original lemma}{original lemmas}
\Crefname{originallemma}{Original lemma}{Original lemmas}
\makeatletter
\newlabel{more-morphisms-lemma-make-good-curves}{{0GK7}{source}{}{sourcepointer-4}{}}
\newlabel{more-morphisms-lemma-make-good-curves@cref}{{[originallemma][0][]0GK7}{[1][0][]source}{}{}{}}
\makeatother
\crefname{originallemma}{original lemma}{original lemmas}
\Crefname{originallemma}{Original lemma}{Original lemmas}
\makeatletter
\newlabel{etale-cohomology-lemma-check-constructible}{{095Q}{source}{}{sourcepointer-5}{}}
\newlabel{etale-cohomology-lemma-check-constructible@cref}{{[originallemma][0][]095Q}{[1][0][]source}{}{}{}}
\makeatother
\crefname{originallemma}{original lemma}{original lemmas}
\Crefname{originallemma}{Original lemma}{Original lemmas}
\makeatletter
\newlabel{lemma-relative-triangle-associated-to-open}{{0GKP}{source}{}{sourcepointer-6}{}}
\newlabel{lemma-relative-triangle-associated-to-open@cref}{{[originallemma][0][]0GKP}{[1][0][]source}{}{}{}}
\makeatother
\crefname{originallemma}{original lemma}{original lemmas}
\Crefname{originallemma}{Original lemma}{Original lemmas}
\makeatletter
\newlabel{lemma-qf-f-shriek-constructible}{{0GK3}{source}{}{sourcepointer-7}{}}
\newlabel{lemma-qf-f-shriek-constructible@cref}{{[originallemma][0][]0GK3}{[1][0][]source}{}{}{}}
\makeatother
\crefname{originallemma}{original lemma}{original lemmas}
\Crefname{originallemma}{Original lemma}{Original lemmas}
\makeatletter
\newlabel{lemma-shriek-and-thickening}{{0GKQ}{source}{}{sourcepointer-8}{}}
\newlabel{lemma-shriek-and-thickening@cref}{{[originallemma][0][]0GKQ}{[1][0][]source}{}{}{}}
\makeatother
\crefname{originallemma}{original lemma}{original lemmas}
\Crefname{originallemma}{Original lemma}{Original lemmas}
\makeatletter
\newlabel{lemma-shriek-composition}{{0F7J}{source}{}{sourcepointer-9}{}}
\newlabel{lemma-shriek-composition@cref}{{[originallemma][0][]0F7J}{[1][0][]source}{}{}{}}
\makeatother
\crefname{originallemma}{original lemma}{original lemmas}
\Crefname{originallemma}{Original lemma}{Original lemmas}
\makeatletter
\newlabel{lemma-proper-smooth-family-curves-modules}{{0GKW}{source}{}{sourcepointer-10}{}}
\newlabel{lemma-proper-smooth-family-curves-modules@cref}{{[originallemma][0][]0GKW}{[1][0][]source}{}{}{}}
\makeatother
\title{Rf\_! of a finite module along a separated affine smooth relative curve over a Noetherian affine base has constructible cohomology sheaves (torsion coefficients)}
\author{The Stacks Project Authors}
\date{stacks-project:more-etale.tex:local-snapshot-20261012; GNU FDL 1.2 or later}
\begin{document}
\maketitle
\setcounter{subsection}{1}
\noindent\textit{Editorial note.} All delivered mathematical text is the author's own, copied from the pinned Stacks Project chapter file more-etale.tex; no statement or proof step was rewritten, reordered or removed. The counted span is the author environment and its own proof environment. References inside the span are handed to the assembler as navigation stubs that carry the canonical tag of the cited tag only; whether each of them is load-bearing is recorded on the stub itself.\par\medskip
\noindent\textit{Allowed standard prerequisites (editorial).}\begin{itemize}
\item Good-curves diagram over a generic point (More on Morphisms, Lemma more-morphisms-lemma-make-good-curves, canonical tag 0GK7). cited at proof lines 4775 and 4784 (part (4)). Supplies the auxiliary diagram over the generic point and the dense open W in the fibres around which the whole reduction chain is organised.
\item Relative triangle associated to an open (Lemma lemma-relative-triangle-associated-to-open, canonical tag 0GKP). cited at proof lines 4788-4795 and 4807-4813. Gives both distinguished triangles used to discard the quasi-finite complement terms Z and Z'.
\item Constructibility for the proper smooth family of curves (Lemma lemma-proper-smooth-family-curves-modules, canonical tag 0GKW). cited at proof lines 4822-4823. The terminal named case to which the induction reduces, combined with finite etale T\_i -> V.
\end{itemize}\medskip
\ifdefined\hypertarget\hypertarget{sourcepointer-1}{}\fi
\noindent\textit{Source reference.} Original lemma 0GK0: sources/stacks-chow/tags.txt, source lines 20670--20670. Cited at proof line 4745: identifies the stalks of R\textasciicircum{}qf\_!M with the compact cohomology H\textasciicircum{}q\_c of the geometric fibre in the base case d = 0.\par\medskip
\ifdefined\hypertarget\hypertarget{sourcepointer-2}{}\fi
\noindent\textit{Source reference.} Original lemma 0GKR: sources/stacks-chow/tags.txt, source lines 20696--20696. Cited at proof line 4746: gives finiteness of H\textasciicircum{}q\_c(X\_s, M) as a Lambda-module in the base case.\par\medskip
\ifdefined\hypertarget\hypertarget{sourcepointer-3}{}\fi
\noindent\textit{Source reference.} Original lemma 0F7L: sources/stacks-chow/tags.txt, source lines 19011--19011. Cited at proof lines 4753-4754: base change compatibility of Rf\_!M lets the induction step be reduced to a dense open U.\par\medskip
\ifdefined\hypertarget\hypertarget{sourcepointer-4}{}\fi
\noindent\textit{Source reference.} Original lemma 0GK7: sources/stacks-chow/tags.txt, source lines 20677--20677. Cited at proof lines 4775 and 4784: supplies the auxiliary diagram over the generic point and its part (4) dense open W in the fibres.\par\medskip
\ifdefined\hypertarget\hypertarget{sourcepointer-5}{}\fi
\noindent\textit{Source reference.} Original lemma 095Q: sources/stacks-chow/tags.txt, source lines 11593--11593. Cited at proof lines 4777-4779: since V -> U is finite surjective, constructibility may be checked after pullback to V.\par\medskip
\ifdefined\hypertarget\hypertarget{sourcepointer-6}{}\fi
\noindent\textit{Source reference.} Original lemma 0GKP: sources/stacks-chow/tags.txt, source lines 20694--20694. Cited at proof lines 4788-4795 and 4807-4813: the two distinguished triangles that discard the quasi-finite complements Z and Z'.\par\medskip
\ifdefined\hypertarget\hypertarget{sourcepointer-7}{}\fi
\noindent\textit{Source reference.} Original lemma 0GK3: sources/stacks-chow/tags.txt, source lines 20673--20673. Cited at proof lines 4796-4798 and 4824-4825: the quasi-finite Z -> V and the finite etale T\_i -> V terms have constructible R\_!M.\par\medskip
\ifdefined\hypertarget\hypertarget{sourcepointer-8}{}\fi
\noindent\textit{Source reference.} Original lemma 0GKQ: sources/stacks-chow/tags.txt, source lines 20695--20695. Cited at proof lines 4801-4804: the thickening nu : nu\textasciicircum{}\{-1\}(W) -> W leaves R\_!M unchanged.\par\medskip
\ifdefined\hypertarget\hypertarget{sourcepointer-9}{}\fi
\noindent\textit{Source reference.} Original lemma 0F7J: sources/stacks-chow/tags.txt, source lines 19009--19009. Cited at proof lines 4816-4821: the second equality factoring R(\textbackslash{}overline\{Y\}\_i -> V)\_!M as R(T\_i -> V)\_! composed with R(\textbackslash{}overline\{Y\}\_i -> T\_i)\_!M.\par\medskip
\ifdefined\hypertarget\hypertarget{sourcepointer-10}{}\fi
\noindent\textit{Source reference.} Original lemma 0GKW: sources/stacks-chow/tags.txt, source lines 20701--20701. Cited at proof lines 4822-4823: the terminal case, constructibility of R(\textbackslash{}overline\{Y\}\_i -> T\_i)\_!M for the proper smooth family of curves.\par\medskip
\begin{lemma}
\label{lemma-constant-shriek-rel-dim-1}
Let $S$ be a Noetherian affine scheme of finite dimension.
Let $f : X \to S$ be a separated, affine, smooth morphism
of relative dimension $1$. Let $\Lambda$ be a Noetherian ring
which is torsion. Let $M$ be a finite $\Lambda$-module.
Then $Rf_!\underline{M}$ has constructible cohomology sheaves.
\end{lemma}

\begin{proof}
We will prove the result by induction on $d = \dim(S)$.

\medskip\noindent
Base case. If $d = 0$, then the only thing to show is that the stalks
of $R^qf_!\underline{M}$ are finite $\Lambda$-modules.
If $\overline{s}$ is a geometric point of $S$, then we have
$(R^qf_!\underline{M})_{\overline{s}} = H^q_c(X_{\overline{s}}, \underline{M})$
by Lemma \ref{lemma-stalk-R-f-shriek}.
This is a finite $\Lambda$-module by Lemma \ref{lemma-finiteness-curves}.

\medskip\noindent
Induction step. It suffices to find a dense open $U \subset S$ such
that $Rf_!\underline{M}|_U$ has constructible cohomology sheaves.
Namely, the restriction of $Rf_!\underline{M}$ to the complement
$S \setminus U$ will have constructible cohomology sheaves by induction
and the fact that formation of $Rf_!\underline{M}$ commutes
with all base change (Lemma \ref{lemma-base-change-shriek}).
In fact, let $\eta \in S$ be a generic point of an irreducible
component of $S$. Then it suffices to find an open neighbourhood $U$
of $\eta$ such that the restriction of $Rf_!\underline{M}$
to $U$ is constructible. This is what we will do in the next paragraph.

\medskip\noindent
Given a generic point $\eta \in S$ we choose a diagram
$$
\xymatrix{
\overline{Y}_1 \amalg \ldots \amalg \overline{Y}_n \ar[rd] &
Y_1 \amalg \ldots \amalg Y_n \ar[r]_-\nu \ar[d] \ar[l]^j &
X_V \ar[r] \ar[d] &
X_U \ar[r] \ar[d] &
X \ar[d]^f \\
& T_1 \amalg \ldots \amalg T_n \ar[r] &
V \ar[r] &
U \ar[r] &
S
}
$$
as in More on Morphisms, Lemma \ref{more-morphisms-lemma-make-good-curves}.
We will show that $Rf_!\underline{M}|_U$ is constructible.
First, since $V \to U$ is finite and surjective, it suffices
to show that the pullback to $V$ is constructible, see
\'Etale Cohomology, Lemma \ref{etale-cohomology-lemma-check-constructible}.
Since formation of $Rf_!$ commutes with base change, we
see that it suffices to show that
$R(X_V \to V)_!\underline{M}$ is constructible.
Let $W \subset X_V$ be the open subscheme given to us by
More on Morphisms, Lemma \ref{more-morphisms-lemma-make-good-curves} part (4).
Let $Z \subset X_V$ be the reduced induced scheme structure on
the complement of $W$ in $X_V$. Then the fibres of $Z \to V$
have dimension $0$ (as $W$ is dense in the fibres) and hence
$Z \to V$ is quasi-finite. From the distinguished triangle
$$
R(W \to V)_!\underline{M} \to
R(X_V \to V)_!\underline{M} \to
R(Z \to V)_!\underline{M} \to
\ldots
$$
of Lemma \ref{lemma-relative-triangle-associated-to-open} and
from Lemma \ref{lemma-qf-f-shriek-constructible}
we conclude that it suffices to show that
$R(W \to V)_!\underline{M}$ has constructible cohomology sheaves.
Next, we have
$$
R(W \to V)_!\underline{M} = R(\nu^{-1}(W) \to V)_!\underline{M}
$$
because the morphism $\nu : \nu^{-1}(W) \to W$ is a thickening
and we may apply Lemma \ref{lemma-shriek-and-thickening}.
Next, we let $Z' \subset \coprod \overline{Y}_i$ denote the
complement of the open $j(\nu^{-1}(W))$. Again $Z' \to V$ is quasi-finite.
Again use the distinguished triangle
$$
R(\nu^{-1}(W) \to V)_!\underline{M} \to
R(\coprod \overline{Y}_i \to V)_!\underline{M} \to
R(Z' \to V)_!\underline{M} \to
\ldots
$$
to conclude that it suffices to prove
$$
R(\coprod \overline{Y}_i \to V)_!\underline{M} =
\bigoplus\nolimits_i R(\overline{Y}_i \to V)_!\underline{M} =
\bigoplus\nolimits_i R(T_i \to V)_!R(\overline{Y}_i \to T_i)_!\underline{M}
$$
has constructible cohomology sheaves (second equality by
Lemma \ref{lemma-shriek-composition}).
The result for $R(\overline{Y}_i \to T_i)_!\underline{M}$
is Lemma \ref{lemma-proper-smooth-family-curves-modules}
and we win because $T_i \to V$ is finite \'etale
and we can apply Lemma \ref{lemma-qf-f-shriek-constructible}.
\end{proof}

\end{document}
